% !TeX root = ../../Book4.tex
% 纲要标题：### 12.2 低阶群上同调与扩张
% 纲要标签：b4:sec:1102
% 纲要位置：计划纲要.md，第 3506 行。

\section{低阶群上同调与扩张}
\label{b4:sec:1102}

低次余链方程可以直接解释为固定元素、交叉同态和群扩张。截面的乘法误差说明 \(2\)-余圈为何出现，也使 Hilbert 第 90 定理具有上同调表述。

% 纲要内容（仅作写作计划，尚非教材正文）：
%
% * H^0、H^1 和 H^2
% * 交叉同态及核为交换群的群扩张
% * 第一卷 Hilbert 第 90 定理的上同调解释
%

\subsection{\texorpdfstring{\(H^0\)、\(H^1\) 与 \(H^2\)}{H0、H1 与 H2}}
\label{b4:subsec:1102:01}

次数零是“不动”的条件，次数一衡量同态方程的扭曲，次数二则衡量截面的乘法缺陷。先把余链方程写清，才能看出这些解释来自同一套微分。

\begin{proposition}\label{b4:prop:group-first-cohomology}
设 \(G\) 是群，\(M\) 是左 \(\mathbb ZG\) 模。则 \(H^0(G,M)=M^G\)，且
\[
H^1(G,M)=
\frac{\{\,c:G\to M\mid c(gh)=c(g)+g c(h)\,\}}
{\{\,g\mapsto ga-a\mid a\in M\,\}}.
\]
分子中的函数称为\term[jiaocha-tongtai]{交叉同态}[crossed homomorphism]，分母中的函数称为主交叉同态。若 \(G\) 在 \(M\) 上平凡作用，则
\[
H^1(G,M)\cong\operatorname{Hom}(G_{\mathrm{ab}},M),
\qquad G_{\mathrm{ab}}=G/[G,G].
\]
\end{proposition}
\begin{proof}
由\eqref{b4:eq:group-cochain-differential}，零次余圈要求 \(ga-a=0\) 对每个 \(g\) 成立，所以得到 \(M^G\)；一次余圈的条件正是所列方程，一次余边是 \(g\mapsto ga-a\)。代入 \(g=h=1\) 可知每个交叉同态均有 \(c(1)=0\)。平凡作用时该方程成为群同态方程，主交叉同态全为零；任何到交换群 \(M\) 的同态都且只通过交换化群分解。
\end{proof}

普通群同态到 \(M\) 的加法群时满足 \(c(gh)=c(g)+c(h)\)，交叉同态则先用 \(g\) 作用于第二个值。这一修正正好来自半直积的乘法。对给定作用，半直积 \(M\rtimes G\) 的乘法为 \((a,g)(b,h)=(a+gb,gh)\)，因此
\[
(c(g),g)(c(h),h)=(c(g)+g c(h),gh).
\]
映射 \(s_c:g\mapsto(c(g),g)\) 为群同态，当且仅当 \(c\) 是交叉同态；它还满足投影与 \(s_c\) 的复合为 \(1_G\)，所以是同态截面。交叉同态方程由此表达的是到半直积的同态条件。

\ProseParagraphBreak
对 \(G=C_2=\langle s\rangle\) 在 \(M=\mathbb Z\) 上的取负作用，交叉同态由任意整数 \(c(s)\) 决定，因为 \(c(s^2)=c(s)-c(s)=0\)；主交叉同态的值为 \(-2a\)。所以 \(H^1(C_2,\mathbb Z_{\mathrm{sgn}})=\mathbb Z/2\mathbb Z\)。这与平凡系数的 \(H^1(C_2,\mathbb Z)=0\) 不同。

\ProseParagraphBreak
次数二的余圈条件和余边公式分别是
\begin{align}
g f(h,\ell)-f(gh,\ell)+f(g,h\ell)-f(g,h)&=0,
\label{b4:eq:group-two-cocycle}\\
(\delta b)(g,h)&=g b(h)-b(gh)+b(g).
\label{b4:eq:group-two-coboundary}
\end{align}
可以让每个二次上同调类都有满足 \(f(1,g)=f(g,1)=0\) 的代表，称为\term[guiyihua-er-yuquan]{归一化 \(2\)-余圈}[normalized 2-cocycle]。事实上，将余圈方程分别用于 \((1,g,h)\)、\((g,1,h)\)，得到 \(f(1,h)=t\)、\(f(g,1)=gt\)，其中 \(t=f(1,1)\)；减去常值一余链 \(b(g)=t\) 的余边即可。两个归一化代表若相差 \(\delta b\)，代入 \((1,g)\) 又得到 \(b(1)=0\)。因而在 \(H^2(G,M)\) 及下面的扩张分类中，可以始终使用归一化 \(2\)-余圈和满足 \(b(1)=0\) 的 \(1\)-余链。

\ProseParagraphBreak
归一化还有群运算上的用途。若集合截面满足 \(s(1)=1_E\)，则 \(s(1)s(g)=s(g)=s(g)s(1)\)，其乘法误差在含单位元的输入处必须为零。反过来，用 \(f\) 定义乘法 \((a,g)(b,h)=(a+gb+f(g,h),gh)\) 时，归一化正好保证 \((0,1)\) 是单位元。未归一化时，例如 \((0,1)(a,g)=(a+f(1,g),g)\)，一般不等于 \((a,g)\)。

\subsection{二次上同调与群扩张}
\label{b4:subsec:1102:02}

\begin{definition}\label{b4:def:extension-fixed-action}
设 \(G\) 是群，\(M\) 是给定的左 \(\mathbb ZG\) 模。一条群的正合列
\[
1\longrightarrow M\xrightarrow{\iota}E\xrightarrow{\pi}G\longrightarrow1
\]
中，若 \(M\) 作为核使用其加法群结构，并且任取 \(\pi(e)=g\) 都有
\(e\iota(m)e^{-1}=\iota(gm)\) 对每个 \(m\in M\) 成立，就称它为实现给定作用的扩张。若 \(\pi\) 存在群同态截面 \(s:G\to E\)，即 \(\pi s=1_G\)，就称扩张分裂。两个这样的扩张的等价，是一个中间群同构与两端的恒等映射组成正合列之间的映射。
\end{definition}

共轭作用不依赖提升的选择。若 \(\pi(e')=\pi(e)\)，则 \(e'e^{-1}\in\ker\pi=\iota(M)\)，所以 \(e'=\iota(a)e\) 对某个 \(a\in M\) 成立。于是
\[
e'\iota(m)e'^{-1}
=\iota(a)\bigl(e\iota(m)e^{-1}\bigr)\iota(a)^{-1}
=e\iota(m)e^{-1}.
\]
最后一个等号使用核的正规性与交换性：括号中的元素仍在 \(\iota(M)\)，因此与 \(\iota(a)\) 交换。两个提升的乘积提升商中两个元素的乘积，故这些共轭自同构也满足作用律。给定作用不要求核位于中心；只有作用平凡时，才得到中心扩张。这里把核和商与既定的 \(M,G\) 识别，也是分类数据的一部分。

\ProseParagraphBreak

这与\cref{b4:thm:ext1-extensions}由 \(\operatorname{Ext}_R^1\) 分类的扩张是不同的问题。那里是左 \(R\) 模之间的短正合列 \(0\to N\to E\to L\to0\)，中间项也有交换的加法及给定的 \(R\) 作用。这里的中间群 \(E\) 一般不交换，商群 \(G\) 也不是该模扩张中的商模，不能直接把这条群扩张当作一次模扩张。群上同调的 Ext 描述是
\[
H^2(G,M)=\operatorname{Ext}_{\mathbb ZG}^{\,2}(\mathbb Z,M),
\]
其中 \(\mathbb Z\) 取平凡作用。次数二有具体的乘法原因：比较集合截面的 \(s(g)s(h)\) 与 \(s(gh)\)，误差要同时输入 \(g,h\) 两个群元素，成为函数 \(f:G^2\to M\)；再比较三个群元素的两种结合方式，就得到这个二变量函数须满足的三变量 \(2\)-余圈方程。它分类的群扩张与上式所计算的二阶模扩张，是两种不同的扩张解释。

\begin{theorem}[群扩张的上同调分类]\label{b4:thm:group-extensions-h2}
设 \(G\) 是群，\(M\) 是给定的左 \(\mathbb ZG\) 模。固定核 \(M\)、商 \(G\) 及给定的作用，以在两端均为恒等映射的扩张同构为等价关系，则实现该作用的扩张的等价类与 \(H^2(G,M)\) 一一对应。零类对应分裂扩张 \(M\rtimes G\)。
\end{theorem}
\begin{proof}
给定扩张，由选择公理选取集合截面 \(s:G\to E\)，使 \(\pi s=1_G\)、\(s(1)=1_E\)。因为 \(s(g)s(h)s(gh)^{-1}\) 的商像为单位元，它在 \(\ker\pi=\iota(M)\) 中；\(\iota\) 单射，故唯一确定函数 \(f:G^2\to M\)，使
\[
s(g)s(h)=\iota(f(g,h))s(gh).
\]
要比较两种结合方式，须把核元素移到所选代表的左侧。由指定的共轭作用，\(s(g)\iota(m)=\iota(gm)s(g)\)，所以
\[
\begin{aligned}
(s(g)s(h))s(\ell)
&=\iota(f(g,h))\iota(f(gh,\ell))s(gh\ell)\\
&=\iota(f(g,h)+f(gh,\ell))s(gh\ell),\\
s(g)(s(h)s(\ell))
&=s(g)\iota(f(h,\ell))s(h\ell)\\
&=\iota(g f(h,\ell)+f(g,h\ell))s(gh\ell).
\end{aligned}
\]
消去共同的右因子，再用 \(\iota\) 单射，就得到\eqref{b4:eq:group-two-cocycle}。作用项 \(g f(h,\ell)\) 正是移过 \(s(g)\) 时产生的；截面在单位元处的条件则给出归一化。任一另选的归一化截面都唯一写为 \(s'(g)=\iota(b(g))s(g)\)，其中 \(b(1)=0\)。代入后有
\[
\begin{aligned}
s'(g)s'(h)
&=\iota\bigl(b(g)+g b(h)+f(g,h)\bigr)s(gh)\\
&=\iota\bigl(f(g,h)+b(g)+g b(h)-b(gh)\bigr)s'(gh).
\end{aligned}
\]
因而 \(f'=f+\delta b\)。改变截面仅改变一个余边，等价扩张传送截面又给出相同函数，故扩张确定良定义的上同调类。

反过来，取归一化 \(2\)-余圈 \(f\)，在集合 \(M\times G\) 上规定
\[
(a,g)(b,h)=(a+g b+f(g,h),gh).
\]
记所得运算为 \(E_f\)。两种三项乘积分别为
\[
\begin{aligned}
\bigl((a,g)(b,h)\bigr)(c,\ell)
&=\bigl(a+gb+(gh)c+f(g,h)+f(gh,\ell),gh\ell\bigr),\\
(a,g)\bigl((b,h)(c,\ell)\bigr)
&=\bigl(a+gb+(gh)c+g f(h,\ell)+f(g,h\ell),gh\ell\bigr).
\end{aligned}
\]
第一坐标相等恰是余圈方程，第二坐标的结合律来自 \(G\)，故运算结合。由归一化，\((0,1)(a,g)=(a,g)=(a,g)(0,1)\)。求 \((a,g)\) 的右逆时，第二坐标须为 \(g^{-1}\)，第一坐标须满足 \(a+gx+f(g,g^{-1})=0\)，所以逆元候选为
\[
\bigl(-g^{-1}a-g^{-1}f(g,g^{-1}),g^{-1}\bigr).
\]
余圈方程在 \((g,g^{-1},g)\) 处给出 \(f(g,g^{-1})=g f(g^{-1},g)\)，故候选逆元左乘 \((a,g)\) 时，第一坐标为
\[
-g^{-1}f(g,g^{-1})+f(g^{-1},g)=0.
\]
它因而也是左逆，\(E_f\) 确为群。归一化又使 \(a\mapsto(a,1)\) 为单同态，投影 \((a,g)\mapsto g\) 为满同态，其核恰为这些 \((a,1)\)。最后，
\[
(a,g)(m,1)=(a+gm,g)=(gm,1)(a,g),
\]
在右侧乘以 \((a,g)^{-1}\) 即得核上的共轭作用为 \(m\mapsto gm\)。所以这是实现给定作用的扩张。

还需说明哪些余圈给出等价的扩张。若 \(f'=f+\delta b\)，且 \(b(1)=0\)，令
\[
T:E_{f'}\longrightarrow E_f,\qquad T(a,g)=(a+b(g),g).
\]
用两套乘法分别计算，得到
\[
\begin{aligned}
T\bigl((a,g)(a',h)\bigr)
&=\bigl(a+ga'+f'(g,h)+b(gh),gh\bigr),\\
T(a,g)T(a',h)
&=\bigl(a+ga'+b(g)+g b(h)+f(g,h),gh\bigr).
\end{aligned}
\]
余边等式使这两式相等。将 \(b\) 换成 \(-b\) 给出逆映射，\(b(1)=0\) 保证 \(T\) 固定核，第二坐标则保证它固定商，所以 \(T\) 是扩张等价。反过来，对任何 \(T:E_{f'}\to E_f\) 的扩张等价，固定商使 \(T(0,g)=(b(g),g)\)，固定核及 \((a,g)=(a,1)(0,g)\) 又迫使 \(T(a,g)=(a+b(g),g)\)，其中 \(b(1)=0\)。保持乘法的同一计算于是给出 \(f'=f+\delta b\)。

原扩张中，给定 \(e\in E\)，先令 \(g=\pi(e)\)，再将 \(es(g)^{-1}\in\iota(M)\) 唯一写为 \(\iota(a)\)，就得到唯一表达 \(e=\iota(a)s(g)\)。映射 \(E_f\to E\)、\((a,g)\mapsto\iota(a)s(g)\) 因而为双射；截面乘法公式使它保持乘法及两端映射。重建扩张的标准截面 \(g\mapsto(0,g)\) 又恢复原来的 \(f\)，故两个构造在等价类上互逆。

最后，\(f=0\) 时乘法就是半直积乘法，且 \(g\mapsto(0,g)\) 是同态截面；扩张分裂时可选同态截面，其 \(f\) 为零。若类为零，写 \(f=\delta b\)，则把截面改为 \(s'(g)=\iota(-b(g))s(g)\) 就使 \(f'=f-\delta b=0\)，也得到同态截面。故零类恰对应分裂扩张。
\end{proof}

上述构造也解释了 \(H^1\)。已知分裂扩张 \(E=M\rtimes G\) 的同态截面对应交叉同态 \(c\)。用核元素 \((a,1)\) 共轭这个截面，因 \((a,1)^{-1}=(-a,1)\)，得到
\[
(a,1)s_c(g)(a,1)^{-1}=(c(g)+a-ga,g).
\]
改变量为 \(a-ga=\delta(-a)(g)\)，属于主交叉同态子群；反之，每个主交叉同态也可由这样的核共轭产生。因此 \(H^1(G,M)\) 分类的是同态截面在核的共轭下的等价类。截面的像是 \(M\) 的补群 \(L\)：它与 \(M\) 的交为单位元，且 \((a,g)=(a-c(g),1)s_c(g)\) 说明 \(ML=E\)。反过来，补群满足 \(\pi|_L:L\to G\) 为同构，其逆给出唯一的同态截面。所以也得到补群的核共轭分类。

\ProseParagraphBreak
这里作为映射的截面与作为子群的补群应当区别。核共轭保持截面的商坐标 \(g\)，而直接用任意 \(x\in E\) 共轭截面后，商坐标一般变成 \(\pi(x)g\pi(x)^{-1}\)。对固定 \(E\) 内的补群，全群共轭却不产生更粗的类别：对任一补群 \(L\)，由 \(E=ML\)，任取 \(x\in E\) 都可写成 \(x=m\ell\)，其中 \(m\in M\)、\(\ell\in L\)。因 \(\ell L\ell^{-1}=L\)，有 \(xLx^{-1}=mLm^{-1}\)。补群的全群共轭与核共轭因而给出相同类别。

\ProseParagraphBreak
\cref{b4:tab:group-cohomology-low-degrees}汇总了前三个次数的解释：\(H^1\) 描述固定分裂扩张中的截面，\(H^2\) 则描述扩张本身，其中零类恰对应分裂扩张。
\begin{table}[htbp]
\centering
\caption{固定群 \(G\) 及左 \(\mathbb ZG\) 模 \(M\) 的低次上同调}
\label{b4:tab:group-cohomology-low-degrees}
\begin{tabularx}{\textwidth}{@{}p{0.08\textwidth}>{\raggedright\arraybackslash}X>{\raggedright\arraybackslash}X@{}}
\toprule
次数 & 代表与等价关系 & 群论意义 \\
\midrule
\(H^0\)
& \(gm=m\) 对所有 \(g\in G\) 成立；无余边
& 不变量子群 \(M^G\) \\
\(H^1\)
& \(c(gh)=c(g)+gc(h)\)；\newline
差为主交叉同态 \(g\mapsto ga-a\)
& \(M\rtimes G\) 中同态截面（等价地，补群）的核共轭类 \\
\(H^2\)
& 归一化 \(2\)-余圈满足\eqref{b4:eq:group-two-cocycle}；\newline
\(f' = f+\delta b\)，其中 \(b(1)=0\)
& 固定核 \(M\)、商 \(G\) 及作用的扩张类，等价保持两端恒等 \\
\bottomrule
\end{tabularx}
\end{table}


\begin{example}\label{b4:exm:cyclic-two-extension}
令 \(G=C_2=\langle s\rangle\)、\(M=\mathbb Z/2\mathbb Z\)，作用平凡。归一化 \(2\)-余圈只剩一个可能非零的值 \(f(s,s)\)，余圈条件对这两种取值都成立。归一化 \(2\)-余边在 \((s,s)\) 处为 \(2b(s)=0\)。故 \(H^2(C_2,\mathbb Z/2\mathbb Z)=\mathbb Z/2\mathbb Z\)：零类给出 \(C_2\times C_2\)，非零类中 \((0,s)^2=(1,1)\)，给出 \(C_4\)。
\end{example}

\subsection{Hilbert 第 90 定理的上同调解释}
\label{b4:subsec:1102:03}

% 跨卷引用目标：b1:thm:hilbert90-multiplicative（20.2.1）、b1:thm:hilbert90-additive（20.2.3）。
设 \(L/K\) 是有限循环 Galois 扩张\footnote{伽罗瓦（Évariste Galois，1811-1832），法国数学家，建立以置换群研究多项式方程可解性的理论。}，\(G=\langle\sigma\rangle\)、\(|G|=m\)。第一卷定理20.2.1和定理20.2.3已经分别证明 Hilbert 第 90 定理的乘法形式与加法形式。这里的 \(L^\times\) 和 \(L\) 都是带 \(G\) 作用的交换群，可以分别作群上同调的系数；前者的群运算使用乘法，后者使用加法。对应的主交叉同态写为
\[
(\delta u)(g)=\frac{g(u)}{u}\quad(u\in L^\times),
\qquad
(\delta v)(g)=g(v)-v\quad(v\in L).
\]
循环关系 \(\sigma^m=1\) 把一次余圈在生成元处的值限制在范数核或迹核中，Hilbert 第 90 定理则使这些余圈成为余边。

\begin{proposition}\label{b4:prop:hilbert90-cohomology}
设 \(L/K\) 是有限循环 Galois 扩张，\(G=\operatorname{Gal}(L/K)\)。对 Galois 自然作用，有
\[
H^1(G,L^\times)=0,\qquad H^1(G,L)=0,
\]
其中第二个系数使用域的加法群。
\end{proposition}
\begin{proof}
选取 \(G\) 的生成元 \(\sigma\)，记 \(|G|=m\)。乘法记号下，一余圈满足 \(c(gh)=c(g)\,g(c(h))\)，且 \(c(1)=1\)。由 \(a=c(\sigma)\) 递推，
\[
c(\sigma^j)=\prod_{i=0}^{j-1}\sigma^i(a)\qquad(0\le j\le m),
\]
其中 \(j=0\) 时取空积 \(1\)。
% 跨卷引用目标：b1:thm:hilbert90-multiplicative（20.2.1）。
与 \(\sigma^m=1\) 相容要求 \(c(\sigma^m)=N_{L/K}(a)=1\)。反过来，范数为一时，指数增加 \(m\) 后新增的一段乘积为 \(\sigma^j(N_{L/K}(a))=1\)，所以值不变；将两段乘积相接又给出余圈方程。因此这也是相容的充分条件。第一卷定理20.2.1给出 \(a=\dfrac{b}{\sigma(b)}\)，其中 \(b\ne0\)。取 \(u=b^{-1}\)，则
\[
\frac{\sigma(u)}{u}
=\frac{\sigma(b^{-1})}{b^{-1}}
=\frac{b}{\sigma(b)}=a.
\]
主交叉同态 \(g\mapsto\dfrac{g(u)}{u}\) 与 \(c\) 在生成元 \(\sigma\) 上取值相同，故由递推在每个 \(g\in G\) 上相同。其值也可由上述乘积相消写为 \(\dfrac{b}{g(b)}\)。两种比值写法只是参数互逆，得到同一余边子群。

% 跨卷引用目标：b1:thm:hilbert90-additive（20.2.3）。
加法情形，交叉同态由 \(a=c(\sigma)\) 决定，递推为
\[
c(\sigma^j)=\sum_{i=0}^{j-1}\sigma^i(a)\qquad(0\le j\le m).
\]
空和为零，与群关系相容的条件是 \(\operatorname{Tr}_{L/K}(a)=0\)；此时增加 \(m\) 项的和为零，连接两段和也给出交叉同态方程。第一卷定理20.2.3给出 \(a=b-\sigma(b)\)。取 \(v=-b\)，就有
\[
\sigma(v)-v=\sigma(-b)-(-b)=b-\sigma(b)=a.
\]
于是生成元处的值及递推给出 \(c(g)=g(v)-v=b-g(b)\)，所以它也是余边。该证明不要求扩张次数在 \(K\) 中可逆。
\end{proof}

对 \(\mathbb C/\mathbb R\)，乘法一次余圈由单位圆上的 \(a\) 决定；Hilbert 第 90 定理断言它总能写成 \(b/\bar b\)。但这不表示所有正次数都消失：下一节将由周期消解算出
\(H^2(C_2,\mathbb C^\times)=\mathbb R^\times/N_{\mathbb C/\mathbb R}(\mathbb C^\times)\cong\mathbb Z/2\mathbb Z\)。

\ProseParagraphBreak
任意有限 Galois 扩张也有乘法 \(H^1\) 消失定理；本小节已经证明的是与第一卷相应的循环版本。更一般的陈述和两种系数的区别，可查 Weibel 的《An Introduction to Homological Algebra》\cite[Hilbert's Theorem 90 6.3.7, Hilbert's Theorem 90 6.4.7, Example 6.4.8]{Weibel1994HomologicalAlgebra}。后续使用这里的结论时仍应核对系数是加法群还是乘法群。
